% Figure 1 of cultivation.tex: the main concepts. Drawn at final size, 8 pt sans-serif text with sans-italic maths,
% in the style of fig_reductions.tex. Requires tikz (arrows.meta, calc, decorations.pathreplacing) and the math
% version figsans declared in the preamble. Colours are CMYK (Okabe-Ito hues). The explanations are in the caption.
\definecolor{tolblue}{cmyk}{1,0.36,0,0.30}      % T states, cat path
\definecolor{harmver}{cmyk}{0,0.56,1,0.165}     % faults and what they spread to
\definecolor{skyblue}{cmyk}{0.63,0.23,0,0.09}   % code, verification
\definecolor{rejgray}{cmyk}{0,0,0,0.62}         % rejected, classical wires
\definecolor{inkgray}{cmyk}{0,0,0,0.75}         % annotation
\providecommand{\figPi}{\text{\fontencoding{OT1}\fontfamily{cmss}\selectfont\char5}}
\begin{tikzpicture}[
  font=\fontsize{8}{9.6}\fontfamily{phv}\selectfont\mathversion{figsans},
  line width=0.5pt,
  >={Stealth[length=1.7mm,width=1.3mm]},
  letter/.style={font=\fontsize{8}{9.6}\fontfamily{phv}\fontseries{b}\selectfont, anchor=north west, inner sep=0pt},
  small/.style={font=\fontsize{7}{8.4}\fontfamily{phv}\selectfont\mathversion{figsans}},
  gate/.style={draw, rounded corners=1pt, fill=black!8},
  tgate/.style={draw=tolblue, rounded corners=1pt, fill=tolblue!12, minimum width=0.36cm, minimum height=0.24cm, inner sep=0pt, font=\fontsize{7}{8}\fontfamily{phv}\selectfont},
  code/.style={draw=skyblue, rounded corners=1pt, fill=skyblue!18},
  cl/.style={double=white, double distance=0.9pt, line width=0.4pt, rejgray},
  brace/.style={decorate, decoration={brace, amplitude=2.2pt}},
  flt/.pic={\fill[harmver, draw=white, line width=0.3pt] (90:0.15) -- (120:0.06) -- (150:0.15) -- (180:0.06) -- (210:0.15) -- (240:0.06) -- (270:0.15) -- (300:0.06) -- (330:0.15) -- (360:0.06) -- (390:0.15) -- (420:0.06) -- cycle;},
  rej/.pic={\draw[rejgray, line width=1.0pt, line cap=round] (-0.09,-0.09) -- (0.09,0.09) (-0.09,0.09) -- (0.09,-0.09);},
  dq/.style={circle, fill=black, inner sep=0pt, minimum size=1.5mm},
  cq/.style={circle, fill=tolblue, draw=white, line width=0.3pt, inner sep=0pt, minimum size=1.6mm},
  sq/.style={rectangle, fill=skyblue, inner sep=0pt, minimum size=1.4mm},
]
% ================================================================ (a) a cultivation circuit and its proxy
\begin{scope}[shift={(0,0)}]
  \node[letter] at (0,3.95) {(a)};
  \node[anchor=west, small, inkgray] at (0.33,3.87) {cultivation};
  % data wires of the color code
  \foreach \y in {2.95,2.65,2.35,2.05} \draw (0.75,\y) -- (5.15,\y);
  \draw[brace] (0.68,1.98) -- (0.68,3.02);
  \node[anchor=east, align=right, small] at (0.62,2.50) {color\\code};
  % injection
  \draw[gate] (0.88,1.92) rectangle (1.30,3.08);
  \node[rotate=90, small] at (1.09,2.50) {inject $|T\rangle$};
  % the check: T layer, measurement of the logical X, T^dagger layer, repeated
  \draw[rounded corners=2pt, dashed, inkgray] (1.45,1.62) rectangle (3.95,3.28);
  \node[anchor=south east, small, inkgray] at (3.95,3.28) {$\times\,n$};
  \foreach \y in {2.95,2.65,2.35,2.05} \node[tgate] at (1.75,\y) {T};
  \draw[gate] (2.10,1.92) rectangle (3.15,3.08);
  \node at (2.625,2.50) {$\bar X$};
  \draw[cl] (2.625,1.92) -- (2.625,1.68);
  \node[anchor=west, small, rejgray] at (2.66,1.72) {$+1$};
  \foreach \y in {2.95,2.65,2.35,2.05} \node[tgate] at (3.55,\y) {T$^\dagger$};
  \pic at (1.98,2.35) {flt};
  % ideal decoder
  \draw[code] (4.10,1.92) rectangle (4.75,3.08);
  \node[small, align=center] at (4.425,2.50) {ideal\\dec.};
  \draw (4.75,2.50) -- (5.15,2.50);
  \node[anchor=west] at (5.15,2.50) {$\tau$};
  % the proxy: every T replaced by S
  \node[tgate, minimum width=0.42cm] (tT) at (1.75,0.75) {T};
  \node[tgate, minimum width=0.42cm, draw=rejgray, fill=black!8] (tS) at (2.75,0.75) {S};
  \draw[->, inkgray] (tT) -- (tS);
  \node[anchor=west, small, inkgray] at (3.10,0.75) {S proxy};
  \draw[inkgray, dotted, line width=0.5pt] (1.75,1.55) -- (tT);
\end{scope}
% ================================================================ (b) one code decides
\begin{scope}[shift={(6.0,0)}]
  \node[letter] at (0,3.95) {(b)};
  \node[anchor=west, small, inkgray] at (0.33,3.87) {one code decides};
  % skeleton: T states as inputs of a Clifford circuit
  \draw[brace] (0.92,2.30) -- (0.92,3.40);
  \node[anchor=east] at (0.86,2.85) {$|T\rangle^{\otimes t}$};
  \foreach \y in {3.30,3.04,2.66,2.40} \draw[tolblue] (1.0,\y) -- (1.30,\y);
  \node[tolblue] at (1.15,2.88) {$\vdots$};
  \draw[gate] (1.30,2.24) rectangle (3.35,3.46);
  \node[small, align=center] at (2.05,2.85) {Clifford\\skeleton};
  \pic at (2.95,3.20) {flt};
  \draw[->, harmver, dotted, line width=0.7pt] (3.11,3.20) -- (3.75,3.20);
  \draw (3.35,2.60) -- (3.75,2.60);
  \node[anchor=west] at (3.75,2.60) {$\tau$};
  % two outcomes: detected, or a Pauli error on the T inputs
  \draw[->, inkgray, line width=0.6pt] (2.32,2.14) -- (1.25,1.80);
  \draw[->, inkgray, line width=0.6pt] (2.32,2.14) -- (3.25,1.30);
  \node[anchor=north] at (0.85,1.80) {$D(F)\neq0$};
  \pic at (1.62,1.58) {rej};
  % the code test
  \draw[brace] (1.82,0.20) -- (1.82,1.12);
  \node[anchor=east] at (1.76,0.66) {$|T\rangle^{\otimes t}$};
  \foreach \y in {1.02,0.78,0.54,0.30} \draw (1.90,\y) -- (3.90,\y);
  \draw[gate, draw=harmver, fill=harmver!12] (2.05,0.17) rectangle (2.80,1.15);
  \node at (2.425,0.66) {$E(F)$};
  \draw[code] (3.00,0.17) rectangle (3.50,1.15);
  \node at (3.25,0.66) {$\figPi_C$};
  \draw[gate] (3.90,0.17) rectangle (4.35,1.15);
  \node at (4.125,0.66) {$V$};
  \draw (4.35,0.66) -- (4.60,0.66);
  \node[anchor=west] at (4.60,0.66) {$\tau$ or $0$ ?};
\end{scope}
% ================================================================ (c) how far one fault spreads
\begin{scope}[shift={(12.0,0)}]
  \node[letter] at (0,3.95) {(c)};
  \node[anchor=west, small, inkgray] at (0.33,3.87) {one fault};
  % GHZ fold: a tree over the data
  \node[small] at (1.40,3.40) {GHZ fold};
  \coordinate (r) at (1.40,2.95);
  \coordinate (n1) at (0.80,2.45);
  \coordinate (n2) at (2.00,2.45);
  \draw[line width=0.7pt] (r) -- (n1); \draw[line width=0.7pt] (r) -- (n2);
  \foreach \x in {1.65,2.00,2.35} \draw[line width=0.7pt] (n2) -- (\x,1.88);
  \foreach \x in {0.45,0.80,1.15} \draw[harmver, line width=1.0pt] (n1) -- (\x,1.88);
  \node[circle, draw, fill=white, inner sep=0pt, minimum size=1.8mm] at (r) {};
  \node[circle, fill=black!45, inner sep=0pt, minimum size=1.5mm] at (n2) {};
  \pic at (n1) {flt};
  \foreach \x in {0.45,0.80,1.15} \node[dq, fill=harmver] at (\x,1.88) {};
  \foreach \x in {1.65,2.00,2.35} \node[dq] at (\x,1.88) {};
  \node[harmver] at (1.40,1.45) {$d=(\delta+1)/2$};
  % verified cat on a path, grown from its middle, end parity measured
  \node[small] at (4.35,3.40) {verified cat};
  \foreach \x in {3.25,3.80,4.35,4.90,5.45} \draw (\x,2.60) -- (\x,1.98);
  \draw[tolblue, line width=1.0pt] (3.25,2.60) -- (5.45,2.60);
  \draw[harmver, line width=1.2pt] (4.62,2.60) -- (5.45,2.60);
  \foreach \x in {3.25,3.80} \node[cq] at (\x,2.60) {};
  \foreach \x in {4.90,5.45} \node[cq, fill=harmver] at (\x,2.60) {};
  \node[circle, draw=tolblue, fill=white, line width=0.7pt, inner sep=0pt, minimum size=2.0mm] at (4.35,2.60) {};
  \pic at (4.62,2.60) {flt};
  \foreach \x in {3.25,3.80,4.35,4.90,5.45} \node[dq] at (\x,1.98) {};
  \draw[skyblue, dashed, line width=0.6pt] (3.25,2.67) .. controls (3.70,2.94) and (5.00,2.94) .. (5.45,2.67);
  \node[sq] at (4.35,2.87) {};
  \node[tolblue] at (4.35,1.45) {$d=\delta$};
  \draw[rejgray, line width=0.3pt] (2.82,1.05) -- (2.82,3.60);
\end{scope}
\end{tikzpicture}
